% !TEX spellcheck = en_US



% ~~~~~~~~~~~~~~~~~~~~~~~~~~~~~~~~~~~~~~~ V
% (NC) 2026 Haluk Bingol
% github.com/halukbingol/LaTeX-Templates
%
% Licensees may copy, distribute, display, and perform the work and make derivative works 
% and remixes based on it only for non-commercial purposes. 
% ~~~~~~~~~~~~~~~~~~~~~~~~~~~~~~~~~~~~~~~ A




% =====================================================================
%  Trees -- Freshman Level: Your First Binary Tree in Java
%  ---------------------------------------------------------------------
%  Built on hblib/Template/hblibTemplate.tex with the libraries in ../../hblib/.
%  Programs and their real outputs are in codeoutput/:
%     codeoutput/<Name>.java  a program (classes without main() are only compiled)
%     codeoutput/<Name>.sh    a terminal session, run in a copy of codeoutput/<Name>/
%     codeoutput/<Name>.txt   what it printed, made by:  bash run_examples.sh
% =====================================================================

\documentclass[aspectratio=169,10pt,compress]{beamer}

% ~~~~~~~~~~~~~~~~~~~~~~~~~~~~~~~~~~~~~~~~~~~~~~~~~~~~~~~~~~~~~~~~~ V hblib
	\usepackage{../../hblib/hblibPresentation} % lib presentation: theme, i/N, \hSection, algorithm2e, ...
	\usepackage{../../hblib/hblibCodeoutput} % lib code + output: \lstcode, \lstterminal, ...
	\usepackage{../../hblib/hblibDatastructures} % lib data structures: \hString, \hStack, \hQueue
% ~~~~~~~~~~~~~~~~~~~~~~~~~~~~~~~~~~~~~~~~~~~~~~~~~~~~~~~~~~~~~~~~~ A hblib


% xcolor, array (and the L{width} column), listings, tikz and algorithm2e come from the hblib libraries
\usepackage[T1]{fontenc}
\usepackage{lmodern}
\usepackage{booktabs}

\newcommand{\win}{\textbf{Windows}}
\newcommand{\mac}{\textbf{macOS}}

% algorithm2e set-up (the package is loaded by hblibPresentation)
\SetAlCapFnt{\small}
\SetAlCapNameFnt{\small}
\SetKwProg{Fn}{function}{}{end}
\SetKw{KwError}{error}

% ---------------------------------------------------------------------
\title{Trees}
\subtitle{Freshman Level: Your First Binary Tree in Java}
\author[Bingol]{Haluk O. Bingol}

\institute[FCIS]{
    Faculty of Computer and Information Sciences\\
    Yeditepe University
}
\date{
	{\tiny \hbTimeStamp}
}

\begin{document}




% ~~~~~~~~~~~~~~~~~~~~~~~~~~~~~~~~~~~~~~ V frm
\begin{frame}[t,fragile]
  \titlepage
\end{frame}
% ~~~~~~~~~~~~~~~~~~~~~~~~~~~~~~~~~~~~~~ A frm




% ~~~~~~~~~~~~~~~~~~~~~~~~~~~~~~~~~~~~~~ V frm
\begin{frame}[t,fragile] \frametitle{Outline}
  \begin{columns}[T]
    \begin{column}{0.48\textwidth}
      \tableofcontents[sections={1-4}, hideallsubsections]
    \end{column}
    \begin{column}{0.48\textwidth}
      \tableofcontents[sections={5-7}, hideallsubsections]
    \end{column}
  \end{columns}
\end{frame}
% ~~~~~~~~~~~~~~~~~~~~~~~~~~~~~~~~~~~~~~ A frm

% ~~~~~~~~~~~~~~~~~~~~~~~~~~~~~~~~~~~~~~ V
\hSection[Setup]{Getting Ready}
% ~~~~~~~~~~~~~~~~~~~~~~~~~~~~~~~~~~~~~~ A




% ~~~~~~~~~~~~~~~~~~~~~~~~~~~~~~~~~~~~~~ V frm
\begin{frame}[t,fragile] \frametitle{What you need}
  \begin{itemize}
    \item A JDK 17 or newer (\texttt{javac} and \texttt{java}), a terminal and an editor or IDE.
    \item One folder, e.g.\ \texttt{trees}, with all \texttt{.java} files of this deck.
  \end{itemize}
  \begin{columns}[T]
    \begin{column}{0.48\textwidth}
      \begin{block}{Windows (Command Prompt)}


% +++++++++++++++++++++++++++++++++++++++ V lst
%: -Lst.
\begin{lstlisting}[style=codesnippet, language={}, basicstyle=\ttfamily\scriptsize]
java -version
cd trees
javac Build.java
java Build
\end{lstlisting}
% +++++++++++++++++++++++++++++++++++++++ A lst


      \end{block}
    \end{column}
    \begin{column}{0.48\textwidth}
      \begin{block}{macOS (Terminal)}


% +++++++++++++++++++++++++++++++++++++++ V lst
%: -Lst.
\begin{lstlisting}[style=codesnippet, language=bash, basicstyle=\ttfamily\scriptsize]
java -version
cd trees
javac Build.java
java Build
\end{lstlisting}
% +++++++++++++++++++++++++++++++++++++++ A lst


      \end{block}
    \end{column}
  \end{columns}

  \lstterminal[basicstyle=\ttfamily\scriptsize]{codeoutput/Check.txt}
\end{frame}
% ~~~~~~~~~~~~~~~~~~~~~~~~~~~~~~~~~~~~~~ A frm

% ~~~~~~~~~~~~~~~~~~~~~~~~~~~~~~~~~~~~~~ V
\hSection[Idea]{The Idea}
% ~~~~~~~~~~~~~~~~~~~~~~~~~~~~~~~~~~~~~~ A




% ~~~~~~~~~~~~~~~~~~~~~~~~~~~~~~~~~~~~~~ V frm
\begin{frame}[t,fragile] \frametitle{A binary tree of numbers}
  \begin{columns}[T]
    \begin{column}{0.44\textwidth}
      \begin{center}
      \begin{tikzpicture}[htree, level distance=0.95cm, level 1/.style={sibling distance=2.4cm}, level 2/.style={sibling distance=1.2cm}, scale=1.0, transform shape]
  \node[hnode, hhot] {8} child {node[hnode] {3} child {node[hnode, hleaf] {1}} child {node[hnode, hleaf] {6}}} child {node[hnode] {10} child[missing] child {node[hnode, hleaf] {14}}};
\end{tikzpicture}
      \end{center}
    \end{column}
    \begin{column}{0.52\textwidth}
      \begin{itemize}
        \item The \textbf{root} is \texttt{8} (green).
        \item Every node has a \textbf{left} and a \textbf{right} child, each may be missing.
        \item \texttt{1}, \texttt{6}, \texttt{14} have no children: \textbf{leaves} (yellow).
        \item \texttt{10} has only a right child.
        \item Below each node hangs a smaller tree: a \textbf{subtree}.
      \end{itemize}
    \end{column}
  \end{columns}
\end{frame}
% ~~~~~~~~~~~~~~~~~~~~~~~~~~~~~~~~~~~~~~ A frm

% ~~~~~~~~~~~~~~~~~~~~~~~~~~~~~~~~~~~~~~ V
\hSection[Build]{Building a Tree in Java}
% ~~~~~~~~~~~~~~~~~~~~~~~~~~~~~~~~~~~~~~ A




% ~~~~~~~~~~~~~~~~~~~~~~~~~~~~~~~~~~~~~~ V frm
\begin{frame}[t,fragile] \frametitle{Step 1: a class for one node}
  \begin{columns}[T]
    \begin{column}{0.52\textwidth}
      \lstcode[basicstyle=\ttfamily\scriptsize]{codeoutput/IntNode.java}{Java}
    \end{column}
    \begin{column}{0.44\textwidth}
      \begin{center}
      \begin{tikzpicture}[box/.style={draw, thick, minimum width=0.9cm, minimum height=0.55cm, font=\small}]
        \node[box, fill=colorTop!50] (v) at (0,0) {8};
        \node[box] (l) at (-0.9,0) {left}; \node[box] (r) at (0.9,0) {right};
        \node[font=\scriptsize] at (-0.9,-0.6) {\texttt{null}}; \node[font=\scriptsize] at (0.9,-0.6) {\texttt{null}};
      \end{tikzpicture}
      \end{center}
      \begin{itemize}\small
        \item \texttt{left} and \texttt{right} are \textbf{references} to other \texttt{IntNode} objects.
        \item A new node has no children: both are \texttt{null}.
      \end{itemize}
    \end{column}
  \end{columns}
\end{frame}
% ~~~~~~~~~~~~~~~~~~~~~~~~~~~~~~~~~~~~~~ A frm




% ~~~~~~~~~~~~~~~~~~~~~~~~~~~~~~~~~~~~~~ V frm
\begin{frame}[t,fragile] \frametitle{Step 2: make the nodes, then link them}
  \lstcoderightstep[0.6][basicstyle=\ttfamily\tiny][basicstyle=\ttfamily\tiny]{codeoutput/Build.java}{Java}{codeoutput/Build.txt}

  \vspace{0.2em}
  \small \texttt{root.left.right} follows two references: from \texttt{8} go left to \texttt{3}, then right to \texttt{6}. \texttt{10} has no left child: \texttt{null}.
\end{frame}
% ~~~~~~~~~~~~~~~~~~~~~~~~~~~~~~~~~~~~~~ A frm

% ~~~~~~~~~~~~~~~~~~~~~~~~~~~~~~~~~~~~~~ V
\hSection[Recursion]{Thinking Recursively}
% ~~~~~~~~~~~~~~~~~~~~~~~~~~~~~~~~~~~~~~ A




% ~~~~~~~~~~~~~~~~~~~~~~~~~~~~~~~~~~~~~~ V frm
\begin{frame}[t,fragile] \frametitle{A tree is made of smaller trees}
  \begin{columns}[T]
    \begin{column}{0.4\textwidth}
      \begin{center}
      \begin{tikzpicture}[htree, level distance=0.95cm, level 1/.style={sibling distance=2.4cm}, level 2/.style={sibling distance=1.2cm}, scale=1.0, transform shape]
  \node[hnode] {8} child {node[hnode, hdone] {3} child {node[hnode, hdone] {1}} child {node[hnode, hdone] {6}}} child {node[hnode] {10} child[missing] child {node[hnode] {14}}};
\end{tikzpicture}
      \end{center}
      \small red: the left subtree of \texttt{8}
    \end{column}
    \begin{column}{0.56\textwidth}
      To count the nodes of a tree:
      \begin{enumerate}\small
        \item if the tree is \textbf{empty} (\texttt{null}): 0 \hfill \emph{base case}
        \item otherwise: 1 (the root) $+$ the nodes of the left subtree $+$ the nodes of the right subtree \hfill \emph{recursive case}
      \end{enumerate}
      \[
      	\mathrm{size}(8) = 1 + \mathrm{size}(3) + \mathrm{size}(10) = 1 + 3 + 2 = 6
      \]
      \small The method trusts itself for the smaller trees. Every call is on a \textbf{smaller} tree, so it ends at the \texttt{null}s.
    \end{column}
  \end{columns}
\end{frame}
% ~~~~~~~~~~~~~~~~~~~~~~~~~~~~~~~~~~~~~~ A frm




% ~~~~~~~~~~~~~~~~~~~~~~~~~~~~~~~~~~~~~~ V frm
\begin{frame}[t,fragile] \frametitle{Size, sum and maximum}
  \begin{columns}[T]
    \begin{column}{0.49\textwidth}
      \lstcode[basicstyle=\ttfamily\tiny, lastline=18]{codeoutput/Recursion.java}{Java}
    \end{column}
    \begin{column}{0.49\textwidth}
      \lstcode[basicstyle=\ttfamily\tiny, firstline=19, firstnumber=19]{codeoutput/Recursion.java}{Java}
    \end{column}
  \end{columns}
\end{frame}
% ~~~~~~~~~~~~~~~~~~~~~~~~~~~~~~~~~~~~~~ A frm




% ~~~~~~~~~~~~~~~~~~~~~~~~~~~~~~~~~~~~~~ V frm
\begin{frame}[t,fragile] \frametitle{Size, sum and maximum: the output}
  \begin{columns}[T]
    \begin{column}{0.4\textwidth}
      \lstoutput[basicstyle=\ttfamily\scriptsize]{codeoutput/Recursion.txt}
    \end{column}
    \begin{column}{0.56\textwidth}
      \begin{itemize}\small
        \item Same pattern three times: handle the empty tree, then combine the root with the answers of the two subtrees.
        \item \texttt{max} has no answer for the empty tree, so it checks the children before calling itself.
        \item $\mathrm{sum} = 8 + 3 + 10 + 1 + 6 + 14 = 42$.
      \end{itemize}
    \end{column}
  \end{columns}
\end{frame}
% ~~~~~~~~~~~~~~~~~~~~~~~~~~~~~~~~~~~~~~ A frm

% ~~~~~~~~~~~~~~~~~~~~~~~~~~~~~~~~~~~~~~ V
\hSection[Traversal]{Visiting Every Node}
% ~~~~~~~~~~~~~~~~~~~~~~~~~~~~~~~~~~~~~~ A




% ~~~~~~~~~~~~~~~~~~~~~~~~~~~~~~~~~~~~~~ V frm
\begin{frame}[t,fragile] \frametitle{Watching preorder work}
  \lstcoderightstep[0.56][basicstyle=\ttfamily\tiny][basicstyle=\ttfamily\scriptsize]{codeoutput/TracePre.java}{Java}{codeoutput/TracePre.txt}

  \vspace{0.2em}
  \small Each level of indentation is one level deeper in the recursion. Every node and every \texttt{null} is reached exactly once.
\end{frame}
% ~~~~~~~~~~~~~~~~~~~~~~~~~~~~~~~~~~~~~~ A frm




% ~~~~~~~~~~~~~~~~~~~~~~~~~~~~~~~~~~~~~~ V frm
\begin{frame}[t,fragile] \frametitle{Three orders}
  \begin{columns}[T]
    \begin{column}{0.36\textwidth}
      \begin{center}
      \begin{tikzpicture}[htree, level distance=0.95cm, level 1/.style={sibling distance=2.2cm}, level 2/.style={sibling distance=1.1cm}, scale=1.0, transform shape]
  \node[hnode] {8} child {node[hnode] {3} child {node[hnode] {1}} child {node[hnode] {6}}} child {node[hnode] {10} child[missing] child {node[hnode] {14}}};
\end{tikzpicture}
      \end{center}
      \uncover<2->{\lstoutput[basicstyle=\ttfamily\scriptsize]{codeoutput/ThreeOrders.txt}}
    \end{column}
    \begin{column}{0.62\textwidth}
      \lstcode[basicstyle=\ttfamily\tiny, lastline=20]{codeoutput/ThreeOrders.java}{Java}
    \end{column}
  \end{columns}
\end{frame}
% ~~~~~~~~~~~~~~~~~~~~~~~~~~~~~~~~~~~~~~ A frm




% ~~~~~~~~~~~~~~~~~~~~~~~~~~~~~~~~~~~~~~ V frm
\begin{frame}[t,fragile] \frametitle{Doing it on paper}
  \begin{columns}[T]
    \begin{column}{0.4\textwidth}
      \begin{center}
      \begin{tikzpicture}[htree, level distance=0.95cm, level 1/.style={sibling distance=2.2cm}, level 2/.style={sibling distance=1.1cm}, scale=1.0, transform shape]
  \node[hnode, horder=1] {8} child {node[hnode, horder=2] {3} child {node[hnode, horder=3] {1}} child {node[hnode, horder=4] {6}}} child {node[hnode, horder=5] {10} child[missing] child {node[hnode, horder=6] {14}}};
\end{tikzpicture}
      \end{center}
      \small numbers: the preorder
    \end{column}
    \begin{column}{0.56\textwidth}
      \begin{itemize}\small
        \item \textbf{preorder}: write the node \emph{when you arrive}, then do the left side, then the right side.
        \item \textbf{inorder}: finish the whole \emph{left} side first, write the node, then the right side.
        \item \textbf{postorder}: finish \emph{both} sides first, write the node last. The root is always last.
        \item Check: the root is first in preorder and last in postorder.
        \item Here inorder gives \texttt{1 3 6 8 10 14}: sorted, because every left subtree is smaller and every right subtree larger
              (a \emph{binary search tree}).
      \end{itemize}
    \end{column}
  \end{columns}
\end{frame}
% ~~~~~~~~~~~~~~~~~~~~~~~~~~~~~~~~~~~~~~ A frm

% ~~~~~~~~~~~~~~~~~~~~~~~~~~~~~~~~~~~~~~ V
\hSection[Mistakes]{Common Mistakes}
% ~~~~~~~~~~~~~~~~~~~~~~~~~~~~~~~~~~~~~~ A




% ~~~~~~~~~~~~~~~~~~~~~~~~~~~~~~~~~~~~~~ V frm
\begin{frame}[t,fragile] \frametitle{Mistake 1: forgetting the \texttt{null} check}
  \lstcoderightstep[0.5][basicstyle=\ttfamily\tiny][basicstyle=\ttfamily\tiny]{codeoutput/Mistake1.java}{Java}{codeoutput/Mistake1.txt}

  \vspace{0.2em}
  \small Java names the problem: \texttt{"t" is null}. A missing child is \texttt{null}, and \texttt{null.left} is impossible. Base case first!
\end{frame}
% ~~~~~~~~~~~~~~~~~~~~~~~~~~~~~~~~~~~~~~ A frm




% ~~~~~~~~~~~~~~~~~~~~~~~~~~~~~~~~~~~~~~ V frm
\begin{frame}[t,fragile] \frametitle{Mistake 2: recursion that does not get smaller}
  \lstcoderightstep[0.54][basicstyle=\ttfamily\tiny][basicstyle=\ttfamily\tiny, lastline=6]{codeoutput/Mistake2.java}{Java}{codeoutput/Mistake2.txt}

  \vspace{0.2em}
  \small \texttt{size(t)} calls \texttt{size(t)} again and again: the call stack fills up. Every recursive call must be on a \textbf{smaller} tree. (The real trace has 1024 lines.)
\end{frame}
% ~~~~~~~~~~~~~~~~~~~~~~~~~~~~~~~~~~~~~~ A frm




% ~~~~~~~~~~~~~~~~~~~~~~~~~~~~~~~~~~~~~~ V frm
\begin{frame}[t,fragile] \frametitle{Mistake 3: \texttt{||} instead of \texttt{\&\&}}
  \lstcoderightstep[0.6][basicstyle=\ttfamily\tiny][basicstyle=\ttfamily\tiny]{codeoutput/Mistake3.java}{Java}{codeoutput/Mistake3.txt}

  \vspace{0.2em}
  \small Tree: our tree plus \texttt{12} and \texttt{16} below \texttt{14}. \texttt{10} has one child, so \texttt{||} calls it a leaf and never looks below it. A leaf has \textbf{no} children: both \texttt{null}.
\end{frame}
% ~~~~~~~~~~~~~~~~~~~~~~~~~~~~~~~~~~~~~~ A frm




% ~~~~~~~~~~~~~~~~~~~~~~~~~~~~~~~~~~~~~~ V frm
\begin{frame}[t,fragile] \frametitle{Mistake 4: linking to the wrong side}
  \lstcoderightstep[0.54][basicstyle=\ttfamily\tiny][basicstyle=\ttfamily\tiny]{codeoutput/Mistake4.java}{Java}{codeoutput/Mistake4.txt}

  \vspace{0.2em}
  \small The second assignment \textbf{replaces} the first: \texttt{3} is lost and \texttt{root.right} is still \texttt{null}. Draw the tree after every link.
\end{frame}
% ~~~~~~~~~~~~~~~~~~~~~~~~~~~~~~~~~~~~~~ A frm




% ~~~~~~~~~~~~~~~~~~~~~~~~~~~~~~~~~~~~~~ V frm
\begin{frame}[t,fragile] \frametitle{Mistakes the compiler finds}
  \lstterminal[basicstyle=\ttfamily\scriptsize]{codeoutput/Typos.txt}

  \vspace{0.1em}
  \small
  \texttt{left} holds an \texttt{IntNode}, not an \texttt{int}: write \texttt{root.left = new IntNode(3)}. Field names must be spelled exactly: \texttt{right}.
\end{frame}
% ~~~~~~~~~~~~~~~~~~~~~~~~~~~~~~~~~~~~~~ A frm

% ~~~~~~~~~~~~~~~~~~~~~~~~~~~~~~~~~~~~~~ V
\hSection[Summary]{Summary}
% ~~~~~~~~~~~~~~~~~~~~~~~~~~~~~~~~~~~~~~ A




% ~~~~~~~~~~~~~~~~~~~~~~~~~~~~~~~~~~~~~~ V frm
\begin{frame}[t,fragile] \frametitle{Cheat sheet}
  \begin{columns}[T]
    \begin{column}{0.56\textwidth}


% +++++++++++++++++++++++++++++++++++++++ V lst
%: -Lst.
\begin{lstlisting}[style=codesnippet, language=Java, basicstyle=\ttfamily\footnotesize]
static int f(IntNode t) {
    if (t == null) return ...;      // empty tree
    ... t.value ...                 // the root
    ... f(t.left) ... f(t.right)    // the subtrees
}
\end{lstlisting}
% +++++++++++++++++++++++++++++++++++++++ A lst


      \small This pattern solves most tree problems: size, sum, height, search, printing.
    \end{column}
    \begin{column}{0.4\textwidth}
      \begin{block}{Do}
        \small check \texttt{null} first; draw the tree; trace small examples; test the empty tree and a single node
      \end{block}
      \begin{block}{Don't}
        \small call the method on the same node; use \texttt{||} for ``leaf''; overwrite a link by accident
      \end{block}
    \end{column}
  \end{columns}
\end{frame}
% ~~~~~~~~~~~~~~~~~~~~~~~~~~~~~~~~~~~~~~ A frm




% ~~~~~~~~~~~~~~~~~~~~~~~~~~~~~~~~~~~~~~ V frm
\begin{frame}[t,fragile] \frametitle{Exercises}
  \begin{enumerate}\small
    \item Draw the tree that \texttt{Build} makes; write its preorder, inorder and postorder by hand, then run \texttt{ThreeOrders}.
    \item Write \texttt{min(IntNode t)} and \texttt{countEven(IntNode t)}.
    \item Write \texttt{height(IntNode t)} (the empty tree has height $-1$). What is the height of the tree on these slides?
    \item Write \texttt{contains(IntNode t, int x)}: is \texttt{x} somewhere in the tree?
    \item Build the tree for the numbers \texttt{5 2 9 1 4} by hand, so that inorder prints them sorted.
    \item Fix \texttt{Mistake1} \dots{} \texttt{Mistake4} and run them again.
  \end{enumerate}
\end{frame}
% ~~~~~~~~~~~~~~~~~~~~~~~~~~~~~~~~~~~~~~ A frm




% ~~~~~~~~~~~~~~~~~~~~~~~~~~~~~~~~~~~~~~ V frm
\begin{frame}[t,fragile]
	\begin{center}
		\vspace{4cm}
		\hRemark{\Huge Questions?}
	\end{center}
\end{frame}
% ~~~~~~~~~~~~~~~~~~~~~~~~~~~~~~~~~~~~~~ A frm



\end{document}
